\section*{Цель работы}
Исследование основных магнитных свойств электротехнической стали, железо-никелевого сплава (пермаллоя) или нанокристаллического сплава на основе железа.

\section*{Теоретические положения}
Ферромагнитные материалы обладают способностью сильно намагничиваться даже в слабых магнитных полях. Это свойство обусловлено наличием доменов — макроскопических областей спонтанной намагниченности. Во внешнем магнитном поле происходит переориентация и смещение границ доменов, что приводит к возникновению результирующего магнитного момента в образце.

Связь между магнитной индукцией $B$ и напряженностью внешнего магнитного поля $H$ является нелинейной и описывается \textit{основной кривой намагничивания} и \textit{петлей гистерезиса}. Основная кривая получается при намагничивании предварительно размагниченного образца. Петля гистерезиса характеризует поведение материала при циклическом перемагничивании.

Ключевой характеристикой является статическая магнитная проницаемость $\mu_\text{ст}$, определяемая как:
\begin{equation}
	\mu_\text{ст} = \frac{B}{\mu_0 H},
	\label{eq:mu_st}
\end{equation}
где $\mu_0 = 4\pi \cdot 10^{-7}$ Гн/м — магнитная постоянная.

При перемагничивании в ферромагнетике возникают потери энергии, состоящие из потерь на гистерезис и потерь на вихревые токи. Суммарные потери определяют площадь петли гистерезиса.

\section*{Описание установки и методики измерений}
Исследования проводятся на установке, схема которой представлена на \cref{fig:scheme}. Установка включает в себя генератор синусоидальных сигналов $G$, милливольтметр $PU$, осциллограф $N$ и испытательный модуль. Исследуемый образец выполнен в виде тороидального сердечника с первичной ($w_1$) и вторичной ($w_2$) обмотками.

\begin{figure}[h!]
	\centering
	\includegraphics[width=0.8\textwidth]{figs/scheme.pdf}
	\caption{Схема установки для исследования магнитных свойств материалов}
	\label{fig:scheme}
\end{figure}

Напряжение $U_X$, подаваемое на горизонтально отклоняющие пластины осциллографа, снимается с резистора $R_T$ и пропорционально току $I_1$ в первичной обмотке. Следовательно, $U_X$ пропорционально напряженности магнитного поля $H$:
\begin{equation}
	H = \frac{U_X \cdot w_1}{2\pi r_\text{ср} R_T},
	\label{eq:H_field}
\end{equation}
где $2\pi r_\text{ср}$ — средняя длина магнитной силовой линии сердечника.

Напряжение $U_Y$, подаваемое на вертикально отклоняющие пластины, снимается с интегрирующей цепи $R_И C_И$. Это напряжение пропорционально магнитной индукции $B$ в образце:
\begin{equation}
	B = \frac{U_Y R_И C_И}{w_2 S},
	\label{eq:B_induction}
\end{equation}
где $S$ — площадь поперечного сечения сердечника.

Одновременная подача напряжений $U_X$ и $U_Y$ на входы осциллографа позволяет наблюдать на его экране зависимость $B(H)$, то есть петлю гистерезиса.

\section*{Порядок выполнения работы}
\begin{enumerate}
	\item \textit{Подготовка к испытаниям.}
	      \begin{itemize}
		      \item Включить питание лабораторного стенда.
		      \item Установить осциллограф в режим $X-Y$.
		      \item Установить частоту генератора $f = 50$ Гц.
		      \item В соответствии с указаниями на стенде установить и зафиксировать масштабы отклонения по осям $m_x$ (В/дел) и $m_y$ (В/дел). Эти значения записать в протокол и не изменять в ходе дальнейших измерений.
		      \item Плавно увеличивая выходное напряжение генератора, получить на экране осциллографа предельную петлю гистерезиса, размах которой по горизонтальной оси составляет 8 делений (от $-4$ до $+4$ дел.). Зарисовать полученную петлю на кальку.
	      \end{itemize}

	\item \textit{Исследование основной кривой намагничивания.}
	      \begin{itemize}
		      \item Уменьшать напряжение генератора таким образом, чтобы координата вершины петли гистерезиса по оси $X$ последовательно принимала значения $4, 3.5, 3, 2.5, 2, 1.5, 1, 0.5, 0$ делений.
		      \item Для каждого значения $X$ измерить соответствующую ему ординату вершины петли по оси $Y$.
		      \item Результаты измерений ($X$, дел. и $Y$, дел.) занести в таблицу протокола (\cref{tbl:mag_curve}).
	      \end{itemize}

	\item \textit{Исследование частотной зависимости магнитных потерь.}
	      \begin{itemize}
		      \item Установить частоту генератора $f=50$ Гц и амплитуду сигнала, соответствующую размаху петли по оси $X$ в 8 делений (насыщение образца).
		      \item Зарисовать петлю гистерезиса на кальку.
		      \item Повторить измерения и зарисовку для частот $200, 400, 600, 800$ Гц, каждый раз устанавливая ту же амплитуду сигнала (размах 8 делений по оси $X$).
		      \item Определить площади полученных петель $S_п$ (например, с помощью миллиметровой бумаги) и занести их в таблицу протокола (\cref{tbl:losses}).
	      \end{itemize}

	\item \textit{Исследование частотной зависимости эффективной магнитной проницаемости.}
	      \begin{itemize}
		      \item Установить переключатель $S$ в положение измерения $U_R$.
		      \item На каждой из частот ($50, 75, 100, 150, 200, 400, 600, 800$ Гц) установить с помощью регулятора амплитуды такое напряжение, чтобы милливольтметр показывал $U_R = 30$ мВ. Это соответствует работе в слабом поле на начальном участке кривой намагничивания.
		      \item Для каждой частоты, не меняя амплитуду, переключить $S$ в положение $U_{BX}$ и измерить входное напряжение.
		      \item Результаты измерений ($f$, Гц и $U_{BX}$, мВ) занести в таблицу протокола (\cref{tbl:mu_eff}).
	      \end{itemize}
\end{enumerate}